% Generated draft: requires mathematical and source review.
\documentclass[11pt]{article}
\usepackage[margin=25mm]{geometry}
\usepackage{fontspec}
\setmainfont{Latin Modern Roman}
\usepackage{amsmath,amssymb,mathtools,unicode-math}
\setmathfont{Latin Modern Math}
\usepackage{enumitem,needspace,xurl,hyperref}
\setlength{\parindent}{0pt}
\setlength{\parskip}{5pt}
\setlength{\emergencystretch}{3em}
\newcommand{\N}{\mathbb{N}}
\newcommand{\Z}{\mathbb{Z}}
\newcommand{\Q}{\mathbb{Q}}
\newcommand{\R}{\mathbb{R}}
\newcommand{\C}{\mathbb{C}}
\newcommand{\F}{\mathbb{F}}
\newcommand{\A}{\mathbb{A}}
\newcommand{\PP}{\mathbb{P}}
\newcommand{\E}{\mathbb{E}}
\newcommand{\eps}{\varepsilon}
\newcommand{\dd}{\,\mathrm{d}}
\begin{document}
\section*{UnsolvedMath ID: 10000082}
\textbf{Source problem code:} catalogue entry
\subsection{Definitions and mathematical statement}
Let \(G_n=(V_n,E_n)\) be finite vertex-transitive graphs with \(N_n=|V_n|\to\infty\). Give edges unit conductance and assume a common finite resistance bound
\[
 \sup_n\sup_{u,v\in V_n}R_{\mathrm{eff}}^{G_n}(u,v)<\infty.
\]
For simple random walk \((X_j)\) on \(G_n\), define the vacant set and largest vacant component size by
\[
 \mathcal V_n(t)=V_n\setminus\{X_0,\ldots,X_{\lfloor t\rfloor}\},\qquad
 L_n(t)=\max\{|C|:C\text{ is a component of }G_n[\mathcal V_n(t)]\},
\]
with \(L_n(t)=0\) if the vacant set is empty.

The source conjectures that, before time \(CN_n\) for a fixed \(C\), the largest vacant component drops from order \(N_n\) to \(o(N_n)\) in a window of width \(o(N_n)\). A strong high-probability interpretation would require times \(t_n^-\leq t_n^+\leq CN_n\), \(t_n^+-t_n^-=o(N_n)\), and a fixed \(c>0\), with
\[
 \Pr(L_n(t_n^-)\geq cN_n)\to1,\qquad
 L_n(t_n^+)/N_n\to0\text{ in probability}.
\]
This display is explicitly an interpretation: the excerpt does not specify deterministic versus random window locations, the probability mode, or uniformity of the lower order constant. It must not be treated as a verified equivalent of every possible sharp-transition formulation.

No degree bound is stated, so the displayed graph class includes dense transitive graphs as well as bounded-degree ones; whether an omitted sparsity hypothesis is needed for the intended linear-time assertion requires source checking. The dependence of \(C\) on the resistance bound or graph sequence is also unstated. The conjecture concerns connected components of the induced vacant graph, not merely the number of unvisited vertices.
\subsection{Short English statement}
Does a random walk on the specified transitive graphs cause the largest unvisited component to pass from macroscopic to sublinear size within a sublinear window before a fixed linear time? The sharp-transition quantifiers and possible missing graph restrictions need verification.
\subsection{Sources}
\begin{itemize}
\item Local UnsolvedMath registry, permanent ID 10000082, source code catalogue entry; Itai Benjamini; no proposal year supplied; updated 2026-08-06.
\item Benjamini, \emph{Random Planar Metrics}, Conjecture 4.1, PDF page 9: \url{https://www.wisdom.weizmann.ac.il/~itai/randomplanar2.pdf#page=9}. Supplied archive: \url{https://arquivo.pt/noFrame/replay/20201231041544id_/http://www.wisdom.weizmann.ac.il/~itai/randomplanar2.pdf}. Registry access date 2026-07-29; extraction checked against rendered PDF, as reported by provenance. Retained after deduplication and source-side review; bounded audit found no complete resolution. Present status and release reuse terms require review.
\item The supplied summary calls this a vacant-set Brussels/dynamical transition and mentions Benjamini--Kozma and expander/transitive-graph work without exact references or a complete theorem.
\end{itemize}
\end{document}
